% Generated draft: requires mathematical and source review.
\documentclass[11pt]{article}
\usepackage[margin=25mm]{geometry}
\usepackage{fontspec}
\setmainfont{Latin Modern Roman}
\usepackage{amsmath,amssymb,mathtools,unicode-math}
\setmathfont{Latin Modern Math}
\usepackage{enumitem,needspace,xurl,hyperref}
\setlength{\parindent}{0pt}
\setlength{\parskip}{5pt}
\setlength{\emergencystretch}{3em}
\newcommand{\N}{\mathbb{N}}
\newcommand{\Z}{\mathbb{Z}}
\newcommand{\Q}{\mathbb{Q}}
\newcommand{\R}{\mathbb{R}}
\newcommand{\C}{\mathbb{C}}
\newcommand{\F}{\mathbb{F}}
\newcommand{\A}{\mathbb{A}}
\newcommand{\PP}{\mathbb{P}}
\newcommand{\E}{\mathbb{E}}
\newcommand{\eps}{\varepsilon}
\newcommand{\dd}{\,\mathrm{d}}
\begin{document}
\section*{UnsolvedMath ID: 30005176}
\textbf{Source problem code:} catalogue entry
\subsection{Definitions and mathematical statement}
Let \(d\in\Z_{>0}\). The supplied question concerns a family \(\pi:X\to S\), with a \(\Q\)-structure implicit in its use of number fields and \(S(k)\). The question does not explicitly specify the admissible class of families or the geometric hypotheses on \(X\), \(S\), and \(\pi\).

For a number field \(k/\Q\) and a point \(s\in S(k)\), write
\[
X_s=X\times_{S,s}\operatorname{Spec}k.
\]
Under the arithmetic-fiber interpretation, \(\operatorname{CH}^i(X_s)\) denotes the group of codimension-\(i\) algebraic cycles on this \(k\)-fiber modulo rational equivalence, and
\[
\operatorname{CH}^i(X_s)_{\mathrm{tors}}
=\{\alpha\in\operatorname{CH}^i(X_s):\text{there exists }n\in\Z_{>0}\text{ with }n\alpha=0\}.
\]
The supplied notation does not settle whether Chow groups of these arithmetic fibers or of their geometric base changes are intended; the display below uses the former interpretation explicitly.

\textbf{Question as supplied, with unresolved family hypotheses.} For each admissible family \(\pi\) and each \(d\in\Z_{>0}\), does there exist an integer \(C=C(\pi,d)>0\) such that
\[
\begin{gathered}
\text{for every number field }k/\Q\text{ with }[k:\Q]\le d,\\
\text{every }s\in S(k),\quad\text{every }i\in\Z_{>0},\quad
\text{and every }\alpha\in\operatorname{CH}^i(X_s)_{\mathrm{tors}},\\
C\alpha=0?
\end{gathered}
\]
Thus a single integer must annihilate all the indicated torsion groups, independently of \(k\), \(s\), and \(i\). This asks for a uniform annihilator, rather than a bound on group cardinalities.

\textbf{Source qualifications.} The surrounding excerpt discusses uniform bounds for cubic surfaces, attributes earlier questions in various forms to Manin, and says that results of Colliot-Th\'el\`ene, Sansuc, and Swinnerton-Dyer imply that the \(2\)-torsion part of codimension-\(2\) Chow groups may be unbounded even in one-parameter families. It also says that uniform boundedness has a negative answer in general. The meaning of ``unbounded'' and the hypotheses of those results are not supplied, so their relationship to the displayed annihilator question is unresolved.

The excerpt separately states a theorem, labelled Theorem 3 and attributed there to reference [1], for a family of smooth projective varieties over a smooth quasi-projective curve over \(\Q\). For a fixed prime \(\ell\) and positive integer \(d\), it asserts a bound, written \(C=C(\pi,d)\), on the cardinality of an \(\ell\)-primary Brauer torsion subgroup for all number fields of degree at most \(d\) and all \(s\in S(k)\) whose fibers satisfy the Tate conjecture for divisors. The group is printed as \(\operatorname{Br}(X_s(k))\), whose intended meaning is unclear. The excerpt additionally asserts an equivalence between finiteness of the \(\ell\)-primary torsion subgroup of \(\operatorname{Br}(X)\) and the \(\ell\)-adic Tate conjecture for divisors, without supplying the necessary conventions or hypotheses. These are contextual source assertions, not hypotheses silently imposed on the Chow-group question or independently verified mathematical claims.
\subsection{Short English statement}
The question asks whether one positive integer, depending only on the family and a bound on number-field degree, kills every torsion Chow class in every positive codimension on every indicated fiber. The family hypotheses and the arithmetic versus geometric interpretation require clarification. The supplied registry labels the problem open; its contextual claims about unbounded \(2\)-torsion and conditional Brauer-group bounds do not resolve these formulation issues.
\subsection{Sources}
\begin{itemize}
\item Local UnsolvedMath registry, supplied complete source record: permanent ID \(30005176\), title ``Uniform Torsion Bounds for Chow Groups in Families,'' exact source problem code \texttt{catalogue entry}. The supplied record labels the problem \texttt{open}, records proposed year 2022, and labels its statement \texttt{corrected\_verified}; its verification note describes repairs to delimiters, notation, grammar, and typography while preserving mathematical content and scope. These are supplied registry labels, not independent mathematical verification.
\item Supplied report citation: \emph{Algebraic Geometry: Moduli Spaces, Birational Geometry and Derived Aspects} (2023), Oberwolfach Reports, pp.\ 1805--1864, DOI 10.4171/owr/2022/32. Supplied source URL: \url{https://doi.org/10.4171/owr/2022/32}. The background describes Oberwolfach Reports as documenting mathematical workshops and their open-problem discussions.
\item Supplied literature source: ``Oberwolfach Report 32/2022,'' \url{https://doi.org/10.4171/owr/2022/32}. Literature assessment checked at \texttt{2026-08-22}: the supplied curation labels the problem open, describes the report's boundedness theorem as conditional on additional hypotheses, recalls possible unbounded \(2\)-torsion in families, and says the stated uniform Chow-torsion bound remains open. This is supplied dated curation, not an independently verified literature assessment; the record itself requests independent verification of time-sensitive status and cited claims.
\item Attribution within the supplied excerpt: Manin for earlier uniform-boundedness questions, and Colliot-Th\'el\`ene, Sansuc, and Swinnerton-Dyer for the cited unboundedness results. No precise references for those results are supplied. The contextual Brauer-group theorem is labelled Theorem 3 ([1]), but the bibliographic identity of [1] is absent. No proposer of the present question is identified.
\end{itemize}
\end{document}
